\chapter{光滑流形与光滑映射}
\label{ch:03-manifolds}
\ChapterPlan{周1.1；梁上册第二版2.1（印刷页15--18）}{流形基础初稿，含例题与解答；未独立数学审阅}



\section{拓扑流形与局部欧氏性}
\begin{definition}[拓扑流形]
固定非负整数$n$。一个无边界的拓扑$n$维流形是Hausdorff、第二可数的拓扑空间$M$，且每个点都有开邻域$U$同胚于$\R^n$的某个开集。这样的同胚称为坐标卡。此定义中的维数$n$固定；维数不变性定理的证明属于另行补充的代数拓扑工具，本节不以直觉代替其证明。
\end{definition}
“局部”不要求整个空间同胚于$\R^n$。例如$S^n$为紧致空间，而$\R^n$在$n\geq1$时不紧致，所以它们不同胚。光滑结构是在拓扑流形上增加的可微数据，不能只靠底层集合确定。

\section{坐标卡、图册与极大光滑结构}
\begin{definition}[光滑流形]
设$M$为Hausdorff、第二可数拓扑空间。一个$n$维坐标卡为$(U,\varphi)$，其中$U\subseteq M$开，$\varphi:U\to\varphi(U)\subseteq\R^n$是到欧氏开集的同胚。覆盖$M$的一族坐标卡称为图册。若任意两卡的过渡映射
\[
\psi\circ\varphi^{-1}:\varphi(U\cap V)\longrightarrow\psi(U\cap V)
\]
均为$C^\infty$映射，则称图册光滑。由所有与该图册光滑兼容的坐标卡组成的极大图册称为光滑结构。配备此结构的$M$称为光滑$n$维流形。
\end{definition}
过渡映射要求两个方向都光滑；上面“任意两卡”同时包含交换两卡后的条件。$U\cap V$为空时没有额外条件。光滑映射之间的复合光滑依赖多元微积分链式法则，其证明在预备章待补。

\begin{theorem}[兼容图册的极大扩充唯一]
每个光滑图册$\mathcal A$包含在唯一的极大光滑图册中。
\end{theorem}
\begin{proof}
令$\mathcal A_{\max}$为所有与$\mathcal A$中每张卡双向光滑兼容的坐标卡的集合。原图册已经互相兼容，故$\mathcal A\subseteq\mathcal A_{\max}$，所以新集合仍覆盖$M$。

取新集合中的两卡$(U,\varphi)$与$(V,\psi)$及$p\in U\cap V$。因原图册覆盖$M$，有$(W,\theta)\in\mathcal A$满足$p\in W$。在$U\cap V\cap W$的坐标像上，
\[
\psi\circ\varphi^{-1}=(\psi\circ\theta^{-1})\circ(\theta\circ\varphi^{-1}).
\]
右侧两映射由兼容性光滑，故左侧在$\varphi(p)$的某邻域光滑。由于$p$任意，过渡映射在整个定义域光滑；交换$\varphi,\psi$同理。因此$\mathcal A_{\max}$是光滑图册。

若某卡与$\mathcal A_{\max}$兼容，它也与其子集$\mathcal A$兼容，所以已在$\mathcal A_{\max}$中；这证明极大性。若$\mathcal B$是包含$\mathcal A$的另一极大光滑图册，其每张卡都兼容于$\mathcal A$，故$\mathcal B\subseteq\mathcal A_{\max}$。极大性迫使等号成立，得唯一性。
\end{proof}



\section{球面、射影空间与乘积流形}
\begin{example}[球面的两张立体投影卡]
设$n\geq1$，$S^n=\{(u,t)\in\R^n\times\R:|u|^2+t^2=1\}$，南北极为$S=(0,-1)$、$N=(0,1)$。在$U_N=S^n\setminus\{N\}$和$U_S=S^n\setminus\{S\}$上定义
\[
\varphi_N(u,t)=\frac{u}{1-t},\qquad
\varphi_S(u,t)=\frac{u}{1+t}.
\]
其逆映射分别为
\begin{align*}
\varphi_N^{-1}(x)&=\left(\frac{2x}{1+|x|^2},\frac{|x|^2-1}{1+|x|^2}\right),\\
\varphi_S^{-1}(x)&=\left(\frac{2x}{1+|x|^2},\frac{1-|x|^2}{1+|x|^2}\right).
\end{align*}
\end{example}
\begin{proof}[坐标卡与兼容性的验证]
由$1\mp t>0$知投影在各自定义域内连续。上式逆映射的分母恒正，故也连续。将分量平方相加得到
\[
\frac{4|x|^2+(|x|^2-1)^2}{(1+|x|^2)^2}=1,
\]
所以逆映射确实取值于球面；其最后一坐标不会等于被删去的极点坐标。代入可验证两个方向复合均为恒等，故投影是到$\R^n$的同胚。

在重叠区域，$x\ne0$。代入得到
\[
\varphi_S\circ\varphi_N^{-1}(x)
=\frac{2x/(1+|x|^2)}{1+(|x|^2-1)/(1+|x|^2)}
=\frac{x}{|x|^2}.
\]
此映射在$\R^n\setminus\{0\}$光滑，并且与自身复合等于恒等，所以逆向过渡也光滑。两卡覆盖球面。球面继承欧氏空间的Hausdorff性；以有理中心、有理半径欧氏开球与球面的交组成可数基，所以球面第二可数。全部流形定义条件均得到验证。
\end{proof}
\begin{figure}[htbp]
\centering
\begin{tikzpicture}[scale=1.75,>=Stealth]
\draw[->,gray] (-1.45,0)--(2.15,0) node[right] {$u$};
\draw[->,gray] (0,-1.25)--(0,1.35) node[above] {$t$};
\draw[thick,blue!65!black] (0,0) circle (1);
\coordinate (N) at (0,1);
\coordinate (P) at (.8,.6);
\coordinate (Q) at (2,0);
\draw[thick,orange!85!black] (N)--(Q);
\draw[dashed] (P)--(.8,0);
\fill (N) circle (.025) node[above left] {$N$};
\fill (P) circle (.025) node[above right] {$P=(u,t)$};
\fill (Q) circle (.025) node[below] {$Q$};
\node[below left] at (0,0) {$O$};
\node[below] at (.8,0) {$u$};
\end{tikzpicture}

\caption{原创示意：立体投影的二维截面。直线$NP$与赤道平面相交于$Q$，其水平坐标为$u/(1-t)$；图仅表示$n=1$的截面关系。}
\label{fig:stereographic}
\end{figure}

\begin{theorem}[开子流形]
若$M$是光滑$n$维流形且$O\subseteq M$开，则$O$带有限制光滑结构，成为光滑$n$维流形。
\end{theorem}
\begin{proof}
子空间$O$继承Hausdorff性与第二可数性。对$M$的卡$(U,\varphi)$，限制卡$(O\cap U,\varphi|_{O\cap U})$的像在$\R^n$中开，因为$O\cap U$在$U$中开且$\varphi$为同胚。这些卡覆盖$O$，过渡映射只是原过渡映射在开子集上的限制，故双向光滑。取其极大扩充即得光滑结构。
\end{proof}
\begin{theorem}[乘积流形]
若$M,N$分别为光滑$m,n$维流形，则$M\times N$自然成为光滑$(m+n)$维流形。
\end{theorem}
\begin{proof}
乘积中不同两点至少一个分量不同，用该因子的Hausdorff分离开集乘上另一因子的全空间，即得不交开邻域。两因子的可数基的所有矩形组成可数基，故乘积第二可数。对卡$(U,\varphi)$与$(V,\psi)$，取乘积卡$\varphi\times\psi:U\times V\to\varphi(U)\times\psi(V)$。其逆是两个逆映射的乘积，故由乘积连续性判别为同胚；像在$\R^{m+n}$中开。两个乘积卡之间的过渡为两个因子过渡映射的乘积，各分量光滑，所以它光滑，反向同理。
\end{proof}
特别地，$T^n=(S^1)^n$是$n$维光滑流形。这里的“自然”指由这些乘积卡生成的光滑结构，具体构造已在证明中给出。

\begin{example}[实射影空间的完整坐标构造]
在$\R^{n+1}\setminus\{0\}$上令$x\sim y$表示$y=\lambda x$，其中$\lambda\ne0$。商空间记为$\mathbb{RP}^n$，其点$[x_0:\cdots:x_n]$代表一条过原点的直线。对$0\leq i\leq n$，令
\[
U_i=\{[x]:x_i\ne0\},\qquad
\varphi_i([x])=(x_0/x_i,\ldots,\widehat{x_i/x_i},\ldots,x_n/x_i).
\]
帽号表示删去该分量。这些表达式不依赖代表元，因为分子分母同时乘以$\lambda$后比值不变。
\end{example}
\begin{proof}[拓扑条件、坐标与过渡函数的验证]
记商映射为$q$。对开集$O\subseteq\R^{n+1}\setminus\{0\}$，
$q^{-1}(q(O))=\bigcup_{\lambda\ne0}\lambda O$开，故由商拓扑定义，$q$是开映射。集合$q^{-1}(U_i)=\{x_i\ne0\}$开，因此$U_i$在商空间中开。其限制$q_i:q^{-1}(U_i)\to U_i$仍是商映射：若$W\subseteq U_i$且$q_i^{-1}(W)$相对开，由$q^{-1}(U_i)$开知该原像在总空间也开，商拓扑定义便推出$W$在总商空间中开，因而在$U_i$中开；反向由$q_i$连续得到。连续比值函数经此商映射下降为连续的$\varphi_i$。反向映射把一个$n$元组送到在第$i$位插入$1$后的等价类，是连续插入映射与$q$的复合。这验证$\varphi_i$为同胚。

补证Hausdorff性。每条直线可由两个相反单位向量表示；对不同的直线$[u]\ne[v]$，取单位代表元，则$\delta=\min\{|u-v|,|u+v|\}>0$。在单位球面上令$A$为$u,-u$的半径$\delta/3$开球邻域之并，$B$为$v,-v$的同类邻域之并。三角不等式表明$A\cap B=\varnothing$。归一化映射$x\mapsto x/|x|$的原像是开且按非零倍数饱和的集合，故分别下降为射影空间的不交开邻域。因此商空间Hausdorff。

每个$U_i$同胚于$\R^n$，具有可数基，而且$U_i$在全空间开。合并有限多个$U_i$的基即得全空间的可数基。最后，在$U_i\cap U_j$上，以第$i$个齐次坐标为$1$，记其余坐标为$y_k$。转到第$j$卡就是将全部齐次坐标除以非零的$y_j$，得到$1/y_j$及各$y_k/y_j$，再删去第$j$位。它们在$y_j\ne0$的开集上光滑，反向换卡同理。故这些卡构成光滑图册。
\end{proof}




\section{光滑映射与微分同胚}
\begin{definition}[光滑映射]
对光滑流形$M,N$，称映射$F:M\to N$光滑，若对每个$p\in M$，存在$p$附近的卡$(U,\varphi)$、$F(p)$附近的卡$(V,\psi)$及$p$的开邻域$U'\subseteq U$，使$F(U')\subseteq V$，且坐标表示$\psi\circ F\circ\varphi^{-1}$在$\varphi(U')$上为$C^\infty$。若$F$双射且$F,F^{-1}$都光滑，称$F$为微分同胚。
\end{definition}
\begin{theorem}[坐标无关性、连续性与复合]
上述定义不依赖选卡。光滑映射连续，光滑映射的复合光滑。
\end{theorem}
\begin{proof}
先在定义给出的$U'$上写出$F=\psi^{-1}\circ(\psi\circ F\circ\varphi^{-1})\circ\varphi$。中间是欧氏空间的光滑映射因而连续，两侧是同胚，故$F|_{U'}$连续。各$U'$覆盖$M$，对$N$中开集$W$，$F^{-1}(W)$与每个$U'$的交开，合并后得$F^{-1}(W)$开，所以$F$连续。
换到其他卡时，在适当缩小的开邻域上新坐标表示为
\[
(\widetilde\psi\circ\psi^{-1})\circ
(\psi\circ F\circ\varphi^{-1})\circ
(\varphi\circ\widetilde\varphi^{-1}).
\]
连续性保证可把定义域限制在这些映射均有定义的开邻域；两端为光滑过渡映射，中间光滑，故由欧氏链式法则复合光滑。复合$G\circ F$时同样缩小邻域，使用同一张中间流形坐标卡，坐标表示就是两个光滑欧氏映射的复合，故光滑。
\end{proof}
这里使用的欧氏链式法则属于预备章仍待补证的分析基础。本节证明解决的是从该分析结果到流形定义的坐标无关性，不宣称已经补齐该分析定理本身的证明。



\section{Hausdorff性与第二可数性的作用}
局部欧氏性不能替代这两项全局条件。
\begin{example}[有两个原点的直线]
取$\R\times\{a,b\}$的两份实线，当$t\ne0$时把$(t,a)$与$(t,b)$识别，保留两个不同原点$0_a,0_b$，赋予商拓扑。每份实线的像都是开集，且自然同胚于$\R$：它的原像是一份完整实线并上另一份去原点的实线；实线中的开集的像也满足同样的开放性检验。因此该空间每点局部像实线。
但$0_a,0_b$的任意开邻域在各自实线中都含一段含零开区间，所以它们都含足够小的非零实数的等价类，必然相交。因此空间不Hausdorff，不是本书约定的流形。
\end{example}
\begin{example}[不可数份实线的不交并]
取不可数指标集$I$和$M=\bigsqcup_{i\in I}\R_i$，规定每份实线开且保持通常拓扑。它Hausdorff且局部欧氏，但非第二可数。若有可数基，则每个非空开集$\R_i$包含至少一个非空基成员$B_i\subseteq\R_i$；不同$i$对应的基成员不同，这要求基不可数，矛盾。
\end{example}


\section{流形基础习题与解答}
第1--3题参照梁书上册第二版2.1的例子及文内留证任务改编（印刷页16--17，PDF第32--33页），不冒称为该章末尾习题的原题号。第4--6题为本工程补充题。
\begin{enumerate}[label=\textbf{3.\arabic*.},ref=3.\arabic*,leftmargin=*]
\Needspace{8\baselineskip}
\item \textbf{圆周的投影坐标。}对$\varepsilon=\pm1$，令$U_{1,\varepsilon}=\{(x,y)\in S^1:\varepsilon x>0\}$并取坐标$y$；类似地令$U_{2,\varepsilon}=\{\varepsilon y>0\}$并取坐标$x$。证明这四张卡组成光滑图册。
\begin{proof}[解答]
各$U$都是圆周相对开集。第一类卡逆映射为$t\mapsto(\varepsilon\sqrt{1-t^2},t)$，$t\in(-1,1)$；第二类为$t\mapsto(t,\varepsilon\sqrt{1-t^2})$。它们连续且与投影互逆。两个第一类卡互不重叠，两个第二类同理。第一类与第二类的过渡映射是在$t>0$或$t<0$的适当开子区间上取$\varepsilon\sqrt{1-t^2}$，其中$|t|<1$，故光滑；反向同形。每个圆周点至少一个坐标非零，所以四卡覆盖圆周；Hausdorff与第二可数性由子空间遗传。
\end{proof}
\Needspace{8\baselineskip}
\item \textbf{两张卡覆盖球面。}验证正文立体投影的过渡映射为$T(x)=x/|x|^2$，并计算其逆。
\begin{proof}[解答]
把$\varphi_N^{-1}$代入$\varphi_S$，分母$1+t=2|x|^2/(1+|x|^2)$，分子$u=2x/(1+|x|^2)$，相除得$T(x)$，定义域为$x\ne0$。因$|T(x)|^2=1/|x|^2$，有$T(T(x))=x$，故$T^{-1}=T$。分母在定义域非零，所以两个方向都光滑。
\end{proof}
\Needspace{8\baselineskip}
\item \textbf{紧致球面不能用一张卡覆盖。}证明$n\geq1$时$S^n$没有定义在整个球面上的坐标卡。
\begin{proof}[解答]
$S^n$为$\R^{n+1}$中的有界闭集，依Heine--Borel定理紧致。若存在全局卡，它将$S^n$同胚到$\R^n$中的非空开集$O$，故$O$紧致又闭。$\R^n$道路连通，因为线段连接任意两点，故$O$既开又闭只能等于$\R^n$。然而$\{B(0,k):k\geq1\}$覆盖$\R^n$却无有限子覆盖，所以$\R^n$不紧致，矛盾。
\end{proof}
\Needspace{8\baselineskip}
\item \textbf{同胚与光滑结构。}在同一实线底空间上分别用全局卡$x(t)=t$和$y(t)=t^3$生成光滑结构。证明两图册不兼容，但两个光滑流形之间存在微分同胚。
\begin{proof}[解答]
$y\circ x^{-1}(s)=s^3$光滑，而$x\circ y^{-1}(s)=\sqrt[3]s$在零点不可微，因为差商为$|s|^{-2/3}$且无有限极限。所以两图册不兼容，底层恒等映射并非双向光滑。但映射$F$从$x$结构实线到$y$结构实线，定义为$F(t)=\sqrt[3]t$，其坐标表示$y\circ F\circ x^{-1}(s)=s$；逆向表示也为恒等，故$F$是微分同胚。不同的底层识别方式不能与“不同微分同胚类型”混为一谈。
\end{proof}
\Needspace{8\baselineskip}
\item \textbf{射影直线的换卡。}计算$\mathbb{RP}^1$两标准卡之间的坐标变换。
\begin{proof}[解答]
在$U_0$上记$t=x_1/x_0$，在$U_1$上记$s=x_0/x_1$。重叠区域要求$x_0x_1\ne0$，故$t,s\ne0$且$s=1/t$。此函数在$\R\setminus\{0\}$光滑，并且逆仍为倒数函数，故双向兼容。
\end{proof}
\Needspace{8\baselineskip}
\item \textbf{第二可数流形的分支数量。}证明每个流形的道路连通分支都是开集，且分支至多可数。
\begin{proof}[解答]
通过一点取坐标卡，再在该点坐标周围取包含于坐标像的开球。此球凸，故逆像道路连通，并是该点的开邻域。若该点属于道路分支$C$，邻域中任一点可由道路连接到该点，再与$C$中的道路拼接，所以邻域包含于$C$。于是$C$开。取可数基，给每个非空$C$选包含于$C$的首个非空基成员；不同分支不交，所选基成员不同，因而道路分支至多可数。道路的拼接在中点连续，可由两段的连续性和它们在端点值相同直接验证。
\end{proof}
\end{enumerate}

